\documentclass[letterpaper,10pt]{article}

% Compiles with pdfLaTeX, including on Overleaf's default settings.
\usepackage{iftex}
\ifPDFTeX
  \usepackage[T1]{fontenc}
  \usepackage[utf8]{inputenc}
  \usepackage[scaled=0.92]{helvet}
\else
  \usepackage{fontspec}
  \setsansfont{texgyreheros}[Extension=.otf, UprightFont=*-regular, BoldFont=*-bold, ItalicFont=*-italic, BoldItalicFont=*-bolditalic, Scale=0.92]
\fi
\renewcommand{\familydefault}{\sfdefault}
\usepackage[letterpaper,margin=0.5in]{geometry}
\usepackage{enumitem}
\usepackage{titlesec}
\usepackage[hidelinks]{hyperref}

% Makes the PDF text map cleanly to Unicode for ATS parsers.
\ifdefined\pdfgentounicode
  \input{glyphtounicode}
  \pdfgentounicode=1
\fi

\pagestyle{empty}
% No line-end hyphenation, so keywords stay whole for ATS parsers.
\hyphenpenalty=10000
\exhyphenpenalty=10000
\emergencystretch=2em
\setlength{\parindent}{0pt}
\setlength{\parskip}{0pt}
\setlist[itemize]{nosep,leftmargin=1.3em,itemsep=0.5pt,topsep=1pt}
\titleformat{\section}{\large\bfseries}{}{0em}{}[\titlerule]
\titlespacing{\section}{0pt}{6pt}{3pt}

\newcommand{\role}[3]{\textbf{#1} \hfill #3\par\textit{#2}\par}

\begin{document}

\begin{center}
{\LARGE\textbf{Ryan McPhedrain}}\par\vspace{4pt}
{\normalsize Machine Learning Scientist and Applied AI Developer}\par\vspace{3pt}
\small
\href{mailto:rmcphedr@gmail.com}{rmcphedr@gmail.com} | 236-512-2798 | Montreal, QC | \href{https://linkedin.com/in/ryanmcphedrain}{linkedin.com/in/ryanmcphedrain} | \href{https://github.com/rmcphedr}{github.com/rmcphedr}
\end{center}

\section{Summary}
Machine learning scientist and applied AI developer with a PhD in Neuroscience and 9+ years of experience turning complex biomedical and behavioral data into actionable insights. Builds reproducible Python pipelines, develops and evaluates statistical and machine learning models, and designs agentic workflows to automate business processes. Brings scientific rigor, practical engineering skills, and clear communication to collaborative research and client projects.

\section{Experience}
\role{Independent AI Consultant}{Montreal, QC}{2026 -- Present}
\begin{itemize}
    \item Built an end-to-end AI-assisted recruiting workflow that consolidates in-house legal openings across multiple sources, ranks opportunities, identifies relevant contacts, and drafts outreach for human review, with an integrated dashboard for monitoring results.
    \item Developed a client-facing platform to support on-demand brand ambassador staffing, centralizing candidate profiles and sourcing workflows while building automated lead discovery and outreach capabilities.
\end{itemize}
\vspace{3pt}

\role{Postdoctoral Researcher}{Mila / McGill University, Montreal, QC}{2024 -- 2026}
\begin{itemize}
    \item Engineered reproducible Python pipelines for large, multi-study human fMRI time-series datasets, integrating temporal filtering, denoising, artifact screening, quality control, feature extraction, statistical analysis, and automated figure generation.
    \item Led data harmonization across 13 studies, 300+ participants, and 500+ scans from teams on 3 continents, resolving differences in acquisition and preprocessing while systematically investigating data-quality and comparability issues.
    \item Implemented a BERT-based semantic harmonization framework to align heterogeneous questionnaire items into shared dimensions, enabling previously incompatible behavioral measures to be analyzed consistently across studies.
    \item Developed multimodal workflows integrating human neuroimaging with rodent transcriptomic data, defining comparable features and validation steps across modalities for cross-species analyses.
    \item Worked directly with collaborating investigators to resolve dataset discrepancies, document analytic decisions, communicate methodological tradeoffs, and translate study-specific requirements into reproducible computational workflows.
\end{itemize}
\vspace{3pt}

\role{Research Consultant (Internship)}{The Decision Lab, Montreal, QC}{2023 (4 months)}
\begin{itemize}
    \item Built and deployed an interactive Python/NLP application over 1,000+ survey responses, translating model outputs into client-facing insights and visual reporting.
    \item Developed and evaluated machine-learning workflows (PCA, K-means, decision trees) for behavioral segmentation, with emphasis on interpretable results and robust comparison across modeling choices.
    \item Published 4 original articles on behavioral science and data-driven decision making, collaborating with interdisciplinary teams on applied research projects.
\end{itemize}
\vspace{3pt}

\role{Graduate Researcher}{McGill University, Montreal, QC}{2016 -- 2023}
\begin{itemize}
    \item Designed and executed experimental neuroscience studies combining in vivo imaging, electrophysiology, molecular biology, and quantitative analysis to investigate mechanisms of developmental neuroplasticity.
    \item Developed Python-based analysis pipelines to automate detection and quantification of neural structures and experimental signals, replacing labor-intensive manual analysis with reproducible computational workflows.
    \item Built statistical analyses linking neural activity, cellular structure, and experimental interventions across longitudinal datasets, with careful attention to signal quality, biological variability, and experimental controls.
    \item Communicated findings through a first-author manuscript and six conference presentations, and collaborated closely with experimental and computational researchers throughout study design, analysis, review, and publication.
\end{itemize}

\section{Projects}
\begin{itemize}
    \item \textbf{Agentic AI Workflow Systems (2026):} Built a personal job-search automation system that consolidates opportunities, evaluates fit against explicit criteria, and tracks applications and follow-ups through a centralized dashboard. Deployed a VPS-hosted Hermes agent integrated with Notion and Google Workspace to coordinate task capture, weekly planning, calendar scheduling, and progress tracking.
\end{itemize}

\section{Skills}
\textbf{Machine Learning and Evaluation:} Python (NumPy, pandas, SciPy, scikit-learn, PyTorch, statsmodels, XGBoost), Classification and Regression, Cross-Validation, Baseline Comparisons, Confounding and Leakage Checks, PCA / ICA, CCA / PLS\par\vspace{2pt}
\textbf{Agentic AI and Engineering:} LLM Integrations, Workflow Automation, Hermes, Cursor / Codex, Structured Evaluation, Logging, Git and GitHub, Docker, Linux, HPC, Reproducible Research Code\par\vspace{2pt}
\textbf{Research and Communication:} Experimental Design, Statistical Modeling, Multi-Study Data Harmonization, Quality Control, Scientific Writing, Data Visualization, Interdisciplinary Collaboration, Client-Facing Reporting

\section{Education}
\role{PhD in Neuroscience (4.0 GPA)}{McGill University, Montreal, QC}{2016 -- 2023}
\role{BSc (Honours) in Molecular Biology and Biochemistry (3.93 GPA)}{Simon Fraser University, Burnaby, BC}{2010 -- 2015}

\section{Awards}
\begin{itemize}
    \item \textbf{1st Place, MedTech Hackathon Challenge (2021):} Designed a computationally efficient wearable device for detecting arrhythmias in performing athletes using built-in physiological signal compression and a lightweight machine-learning model suitable for edge inference.
    \item \textbf{2nd Place, MedTech Hackathon Challenge (2023):} Developed a clinical decision-support system for emergency-department patient flow using predictive modeling to improve inpatient bed-allocation decisions.
\end{itemize}

\end{document}
